\documentclass[10pt,a4paper]{article}

\usepackage[a4paper,margin=1.8cm]{geometry}
\usepackage[T1]{fontenc}
\usepackage{lmodern}
\usepackage{microtype}
\usepackage{graphicx}
\usepackage{booktabs,tabularx,longtable,array,multirow}
\usepackage{amsmath}
\usepackage{siunitx}
\usepackage{caption,subcaption,float}
\usepackage{enumitem}
\usepackage{titlesec}
\usepackage{xcolor}
\usepackage[hidelinks]{hyperref}
\usepackage[nameinlink,noabbrev]{cleveref}
\usepackage{fancyhdr}
\usepackage{url}

\graphicspath{{../outputs/figures/}}
\setlength{\parindent}{0pt}
\setlength{\parskip}{3pt}
\setlist{nosep,leftmargin=1.3em}
\captionsetup{font=small,labelfont=bf,skip=3pt}
\titlespacing*{\section}{0pt}{7pt}{3pt}
\titlespacing*{\subsection}{0pt}{5pt}{2pt}
\titleformat{\section}{\large\bfseries}{\thesection}{0.55em}{}
\titleformat{\subsection}{\normalsize\bfseries}{\thesubsection}{0.45em}{}
\renewcommand{\arraystretch}{1.12}
\sisetup{round-mode=places,round-precision=3}

\pagestyle{fancy}
\fancyhf{}
\fancyhead[L]{FRAM Assignment -- Biocon Limited}
\fancyhead[R]{Yash Agarwal}
\fancyfoot[C]{\thepage}
\setlength{\headheight}{13pt}

\newcommand{\bp}{\,\text{bp}}
\newcommand{\sigE}{\sigma_E}
\newcommand{\sigV}{\sigma_V}
\newcommand{\Rupee}{INR\,}

\begin{document}

% ---------- Unnumbered front matter ----------
\pagenumbering{gobble}
\begin{titlepage}
\centering
\vspace*{0.7cm}
{\Large\bfseries BIRLA INSTITUTE OF TECHNOLOGY \& SCIENCE, PILANI\par}
\vspace{0.2cm}
{\large\bfseries HYDERABAD CAMPUS\par}
\vspace{1.4cm}
{\Large\bfseries FINANCIAL RISK ANALYTICS AND MANAGEMENT\par}
\vspace{0.25cm}
{\large FIN F414\par}
\vspace{1.4cm}
{\large\bfseries ASSIGNMENT\par}
\vspace{0.5cm}
{\LARGE\bfseries Four Routes to Default Probability\par}
\vspace{0.35cm}
{\Huge\bfseries BIOCON LIMITED\par}
\vspace{1.6cm}

\begin{tabularx}{0.92\textwidth}{Xr}
\textit{Submitted by} & \textit{Submitted to}\\[2pt]
\textbf{Yash Agarwal} & \textbf{Prof. A. C. V. Subrahmanyam}\\
ID: 2024A2PS1030H & \\
\end{tabularx}
\vfill
{\large\textbf{Date of submission}\par}
\vspace{0.15cm}
{\large 18 November 2026\par}
\end{titlepage}
\pagestyle{empty}

\section*{Abstract}
This report estimates Biocon Limited's one-year probability of default through four complementary routes at 30 September 2026. First, an independently implemented Black--Scholes--Merton solver extracts at-the-money option-implied equity volatility and feeds it into the Merton structural model. Second, rolling Student-$t$ GARCH(1,1) and Random Forest forecasts provide alternative one-year volatility inputs. Third, Biocon's CRISIL rating and the associated historical default study are reverse-calibrated into Merton-consistent volatility. Fourth, Altman accounting scores are mapped to a rating-equivalent PD and reverse-calibrated in the same manner. The market routes produce reported PDs at the mandatory 3 bp floor, although their distances to default differ materially. The CRISIL route also reports 3 bp after flooring, while Altman EM gives 7 bp. Random Forest does not outperform GARCH out of sample. The comparison shows that market, rating and accounting measures capture different risk horizons; Student-$t$ GARCH-Merton DD is therefore recommended as the primary early-warning indicator, supported by rating and accounting measures.

\clearpage
\tableofcontents
\clearpage

% ---------- Numbered technical report ----------
\pagenumbering{arabic}
\setcounter{page}{1}
\pagestyle{fancy}

\section{Executive Summary}
Biocon Limited's one-year default risk is estimated through point-in-time market measures, volatility forecasts, a through-the-cycle agency rating and an accounting score. \Cref{tab:summary} reports the common 30 September 2026 comparison. The first three routes estimate volatility and then structural PD; CRISIL and Altman begin with an external PD and reverse-engineer the Merton-consistent volatility.

\begin{table}[H]
\centering
\caption{Comparison of one-year default-risk routes at 30 September 2026}
\label{tab:summary}
\scriptsize
\setlength{\tabcolsep}{3.2pt}
\begin{tabularx}{\textwidth}{l l r r r r X}
\toprule
Route & Volatility or PD input & $\sigE$ & $\sigV$ & DD & PD & Perspective\\
\midrule
ATM implied volatility & Interpolated call/put ATM IV & 24.915\% & 21.653\% & 9.281 & 3 bp & Risk-neutral, point-in-time\\
Student-$t$ GARCH & 252-day expected average variance & 32.690\% & 28.410\% & 7.014 & 3 bp & Risk-neutral forecast\\
Random Forest & 252-day realised-volatility forecast & 31.480\% & 27.359\% & 7.295 & 3 bp & Risk-neutral forecast\\
CRISIL AA+ & AA category PD after floor & 63.124\% & 54.862\% & 3.432 & 3 bp & Through-the-cycle reverse calibration\\
Altman EM & A-equivalent category PD & 67.098\% & 58.319\% & 3.195 & 7 bp & Accounting reverse calibration\\
\bottomrule
\end{tabularx}
\end{table}

The reported 3 bp observations are not economically identical: the floor conceals substantial differences in raw PD and DD. ATM IV gives the lowest volatility and largest DD, while GARCH and Random Forest imply moderately greater risk. Random Forest fails to beat GARCH on the untouched 252-day test sample. CRISIL and Altman require much higher reverse-calibrated volatility because they start from historical category PDs rather than current option or equity dynamics. Altman EM is the most conservative route. For continuing surveillance, Student-$t$ GARCH-Merton DD is preferred because it is timely, responsive and empirically stronger than Random Forest; CRISIL and Altman should be retained as slower-moving challenge indicators. The conclusion is not that default is impossible, but that Biocon's market-implied asset cushion is large relative to the assignment default point.

\newpage
\section{Data, Conventions and Assumptions}
The analysis uses 12 last-trading-day snapshots from October 2025 to September 2026. Only information public by each date is used. Adjusted prices and returns supplied by the user are treated as already corrected for corporate actions and are not adjusted again; contemporaneous unadjusted close is used for market capitalisation. Exact shares outstanding change on verified allotment dates. Consolidated debt is taken from the latest publicly released balance sheet and carried forward without forecasting until the next publication.

RBI 91-day and 364-day auction cut-off yields are used for BSM and one-year Merton respectively. Published percentage yields are converted to continuous rates as $r=\ln(1+R/100)$. Trailing-twelve-month cash dividend per share is divided by spot. Market equity is $E=S\times\text{shares}$ and the KMV default point is
\[
D=\text{short-term borrowings}+0.5\times\text{long-term borrowings}.
\]
The main PD is risk-neutral; physical drift is reported as a robustness result. Raw PD below 0.03\% is reported at the mandatory 3 bp floor.

\subsection{Model assumptions and applicability}
BSM assumes geometric Brownian motion, normally distributed log returns, constant volatility and rates, known continuous dividend yield, frictionless continuous trading, unrestricted short selling, no persistent arbitrage, European exercise and synchronous fair prices. These conditions hold only approximately for Biocon: the volatility smile, heavy-tailed returns, discrete dividends, bid--ask effects and time-varying volatility limit literal interpretation. If they fail, a single IV cannot price all strikes and may incorporate liquidity or timing noise.

Merton assumes firm assets follow GBM with constant $\sigV$, the firm has one zero-coupon debt claim, default occurs only at $T$ when $V<D$, and equity is a European call on assets. Biocon has multiple liabilities, refinancing exposure, taxes, dividends and possible distress before $T$. Consequently, the default point compresses the debt structure and PD is a model-based risk indicator, not an exact forecast. A complete assumption-by-assumption assessment is provided in \Cref{app:assumptions}.

\subsection{Data audit}
The raw F\&O workbook contains 1,541 rows: 1,505 options and 36 futures. Sequential option cleaning removes 859 wrong-expiry rows, 199 zero-volume rows, 143 contracts outside $0.90\leq K/S\leq1.10$, and 2 lower-bound violations, leaving 302 BSM-ready contracts. Upper-bound checks are retained as diagnostics. The main source and timing schedule is listed in \Cref{app:data}.

\newpage
\section{Part A: Implied Volatility and Merton PD}
\subsection{Option cleaning and implied volatility}
For a dividend-paying stock, the BSM call price is $c=Se^{-qT}N(d_1)-Ke^{-rT}N(d_2)$, with
\[
d_1=\frac{\ln(S/K)+(r-q+\tfrac12\sigma^2)T}{\sigma\sqrt{T}},\qquad d_2=d_1-\sigma\sqrt{T}.
\]
The near-month expiry is selected when more than seven calendar days remain; otherwise the next expiry is used. Zero-volume contracts, moneyness outside the assignment band and lower-bound violations are removed. The own IV solver starts Newton--Raphson at 30\%, requires absolute price error below INR 0.0001 and switches to bisection on $[0.01,3]$ when necessary. All 302 contracts converge. The maximum own-solver difference from check-only \texttt{py\_vollib} is approximately $3.37\times10^{-6}$.

ATM IV is interpolated separately for calls and puts between strikes surrounding $K/S=1$, then averaged. The 30 September 2026 ATM IV is 24.915\%. August 2026 remains a disclosed data-quality concern because call and put ATM IV differ materially.

\begin{figure}[H]
\centering
\begin{subfigure}{0.49\textwidth}\includegraphics[width=\linewidth]{iv_smile_2026_09_30.png}\caption{Call and put IV smile}\end{subfigure}\hfill
\begin{subfigure}{0.49\textwidth}\includegraphics[width=\linewidth]{atm_iv_monthly.png}\caption{Twelve-date ATM IV}\end{subfigure}
\caption{Biocon option-implied volatility results}\label{fig:iv}
\end{figure}

\subsection{Structural calibration}
Merton equity satisfies
\[
E=VN(d_1)-De^{-rT}N(d_2),\qquad \sigE E=N(d_1)\sigV V.
\]
The two equations are solved jointly using SciPy \texttt{fsolve}, starting at $V_0=E+D$ and $\sigma_{V,0}=\sigE E/(E+D)$. Each solution is substituted back into both equations and independently checked with the assignment-listed \texttt{merton} package. Distance to default is
\[
DD=\frac{\ln(V/D)+(\mu-\tfrac12\sigV^2)T}{\sigV\sqrt{T}},\qquad PD=N(-DD).
\]
At the main date, $E=\Rupee603{,}475.738$ million, $D=\Rupee96{,}530$ million, $V=\Rupee694{,}387.569$ million, $\sigV=21.653\%$, risk-neutral $DD=9.281$ and physical $DD=9.397$. Both raw PDs are below the floor. The high DD reflects the asset-value cushion relative to the default point, not an absence of borrowings; observed $D/V$ is approximately 12--18\% across the 12 dates.

\newpage
\begin{figure}[H]
\centering
\begin{subfigure}{0.49\textwidth}\includegraphics[width=\linewidth]{merton_dd_over_time.png}\caption{DD over the 12 dates}\end{subfigure}\hfill
\begin{subfigure}{0.49\textwidth}\includegraphics[width=\linewidth]{merton_dd_vs_sigma_e_2026_09_30.png}\caption{DD sensitivity to $\sigE$}\end{subfigure}
\caption{Merton structural-risk results}\label{fig:merton}
\end{figure}

\begin{table}[H]
\centering
\caption{Main-date Merton calibration under ATM IV}
\small
\begin{tabular}{lrrrrrrr}
\toprule
$E$ (INR mn) & $D$ (INR mn) & $D/V$ & $\sigE$ & $\sigV$ & $DD_{RN}$ & $DD_{P}$ & PD\\
\midrule
603,475.738 & 96,530.000 & 13.901\% & 24.915\% & 21.653\% & 9.281 & 9.397 & 3 bp\\
\bottomrule
\end{tabular}
\end{table}

Increasing equity volatility monotonically reduces DD. The sensitivity figure therefore provides more economic information than the clipped PD: once raw PD lies below the floor, DD continues to distinguish low-risk from very-low-risk structural states. The physical DD is slightly higher because the historical mean return exceeds the risk-free drift at the main date. Full 12-date results and validation tolerances are reported in \Cref{app:results}.

\section{Part B: GARCH and Random Forest}
\subsection{Rolling Student-$t$ GARCH(1,1)}
Daily adjusted log returns are modelled as $\sigma_t^2=\omega+\alpha u_{t-1}^2+\beta\sigma_{t-1}^2$. Normal and Student-$t$ errors are estimated on the previous 1,000 trading days at each valuation date. Student-$t$ has lower BIC on all 12 dates and is selected because the fitted degrees of freedom capture heavy tails. The annualised input is based on the average variance expected over the next 252 trading days using the assignment's Hull formula.

At 30 September 2026, GARCH gives $\sigE=32.690\%$, $\sigV=28.410\%$, risk-neutral $DD=7.014$ and reported PD of 3 bp. ARCH-LM rejects homoskedasticity in raw returns, supporting GARCH. Ljung--Box tests on squared standardised residuals do not reject at lags 10 or 20, although short-lag correlation remains in standardised residuals and is a limitation.

\newpage
\begin{figure}[H]
\centering
\begin{subfigure}{0.49\textwidth}\includegraphics[width=\linewidth]{garch_sigma_e_over_time.png}\caption{Student-$t$ GARCH volatility}\end{subfigure}\hfill
\begin{subfigure}{0.49\textwidth}\includegraphics[width=\linewidth]{garch_merton_dd_over_time.png}\caption{GARCH-Merton DD}\end{subfigure}
\caption{Rolling GARCH results}\label{fig:garch}
\end{figure}

\subsection{Random Forest volatility forecasting}
Random Forest is chosen for nonlinear interactions, robustness through ensemble averaging and directly reportable importance in a modest sample. Targets are log annualised realised volatility over 21 and 252 days. Features known at $t$ include realised volatility over 5, 21 and 63 days; returns over 1, 5 and 21 days; filtered GARCH volatility; India VIX; Nifty 50 21-day return; and a results-month indicator. Observations are split chronologically 60\%/20\%/20\%, with overlapping forward labels purged and no shuffling. Hyperparameters are selected on validation QLIKE; the untouched test set is evaluated once.

\begin{table}[H]
\centering
\caption{Out-of-sample volatility forecasting performance}
\label{tab:ml}
\small
\begin{tabular}{clrrrr}
\toprule
Horizon & Model & RMSE & MAE & QLIKE & $R^2_{OOS}$\\
\midrule
21-day & Random Forest & 0.08685 & 0.06954 & 0.15926 & -0.00135\\
21-day & GARCH & 0.08679 & 0.07227 & 0.12841 & 0.00000\\
252-day & Random Forest & 0.02430 & 0.01811 & 0.01076 & -0.49214\\
252-day & GARCH & 0.01990 & 0.01520 & 0.00601 & 0.00000\\
\bottomrule
\end{tabular}
\end{table}

Random Forest does not beat GARCH on the required one-year horizon. Its 30 September forecast gives $\sigE=31.480\%$, $\sigV=27.359\%$, $DD=7.295$ and reported PD of 3 bp. Equal weighting remains the locked specification; exponential recency weighting is reported only as a post-model sensitivity.

\begin{figure}[H]
\centering
\begin{subfigure}{0.49\textwidth}\includegraphics[width=\linewidth]{ml_feature_importance_252.png}\caption{252-day feature importance}\end{subfigure}\hfill
\begin{subfigure}{0.49\textwidth}\includegraphics[width=\linewidth]{ml_vs_garch_merton_dd.png}\caption{GARCH and RF Merton DD}\end{subfigure}
\caption{Random Forest interpretation and structural-risk comparison}\label{fig:ml}
\end{figure}

\newpage
\section{Part C: Rating-Implied PD and Volatility}
Biocon's CRISIL AA+/Stable rating was reaffirmed on 16 December 2025 and remained the relevant long-term rating at the main date. The FY2026 default study reports one-year CDRs of 0.00\% for AAA, 0.02\% for AA and 0.07\% for A. AA+ maps to the AA category. After applying the 3 bp floor, the point and lower edge are 3 bp; the upper geometric edge is $\sqrt{3\times7}=4.5826$ bp.

The target PD is converted to $DD=-N^{-1}(PD)$. Holding 30 September 2026 $E$, $D$, $r$ and $T$ fixed, the Merton equity equation and target-DD equation are solved jointly for $(V,\sigV)$; $\sigE$ follows from the Merton volatility identity. This is an inverse calibration: PD is the input and volatility is the implied output.

\begin{table}[H]
\centering
\caption{Risk-neutral rating reverse calibration}
\small
\begin{tabular}{lrrrr}
\toprule
Case & PD & DD & $\sigV$ & $\sigE$\\
\midrule
Lower edge & 3.000 bp & 3.432 & 54.862\% & 63.124\%\\
AA point & 3.000 bp & 3.432 & 54.862\% & 63.124\%\\
Upper edge & 4.583 bp & 3.315 & 56.520\% & 65.024\%\\
\bottomrule
\end{tabular}
\end{table}

\begin{figure}[H]
\centering
\includegraphics[width=0.72\textwidth]{rating_implied_volatility_band.png}
\caption{Biocon market volatility versus the CRISIL-implied band}\label{fig:rating}
\end{figure}

ATM IV and GARCH lie below the 63.124--65.024\% rating-consistent band. This does not prove either measure is stale: option and GARCH measures are point-in-time, whereas ratings and category default studies are through-the-cycle and incorporate qualitative judgement. The band is therefore a credit-consistency benchmark, not an observable market forecast.

\newpage
\section{Part D: Altman Accounting Route}
The original $Z$, private-firm $Z''$ and emerging-market score are
\[
Z=1.2X_1+1.4X_2+3.3X_3+0.6X_4+X_5,\quad
Z''=6.56X_1+3.26X_2+6.72X_3+1.05X_4,\quad EM=3.25+Z''.
\]
Working capital is current assets less current liabilities; $X_2$ uses consolidated Other equity as reserves and surplus according to the assignment; EBIT equals PBT plus finance costs. Original $Z$ uses market equity, while $Z''$ uses book equity. EM is preferred because Biocon is an emerging-market company.

\begin{table}[H]
\centering
\caption{FY2026 Altman results}
\small
\begin{tabular}{lrrl}
\toprule
Model & Score & Zone & Credit interpretation\\
\midrule
Original $Z$ & 2.428 & Grey & Indeterminate under original listed-manufacturer cut-offs\\
$Z''$ & 3.597 & Safe & Low distress signal\\
Emerging Market & 6.847 & Safe & A-equivalent under Altman mapping\\
\bottomrule
\end{tabular}
\end{table}

\begin{figure}[H]
\centering
\begin{subfigure}{0.32\textwidth}\includegraphics[width=\linewidth]{altman_z_score.png}\caption{Original $Z$}\end{subfigure}\hfill
\begin{subfigure}{0.32\textwidth}\includegraphics[width=\linewidth]{altman_z_double_prime.png}\caption{$Z''$}\end{subfigure}\hfill
\begin{subfigure}{0.32\textwidth}\includegraphics[width=\linewidth]{altman_em_score.png}\caption{EM score}\end{subfigure}
\caption{Biocon Altman scores, FY2022--FY2026}\label{fig:altman}
\end{figure}

FY2026 EM maps to A, four notches below the observed AA+ agency rating. The CRISIL A-category one-year PD is 7 bp, giving $DD=3.195$. Reverse calibration produces $\sigV=58.319\%$ and $\sigE=67.098\%$. The gap arises because Altman is a mechanical transformation of a small set of historical accounting ratios, whereas CRISIL also evaluates business profile, liquidity, management, industry conditions and forward-looking factors. The A-equivalent mapping is not Biocon's actual rating.

\newpage
\section{Part E: Integrated Comparison}
\begin{figure}[H]
\centering
\begin{subfigure}{0.49\textwidth}\includegraphics[width=\linewidth]{part_e_volatility_comparison.png}\caption{Equity volatility and CRISIL band}\end{subfigure}\hfill
\begin{subfigure}{0.49\textwidth}\includegraphics[width=\linewidth]{part_e_pd_comparison.png}\caption{Reported PD on log scale}\end{subfigure}
\caption{Cross-route comparison at 30 September 2026}\label{fig:parte}
\end{figure}

The volatility ordering is ATM IV $<$ Random Forest $<$ GARCH $<$ CRISIL $<$ Altman. Relative to ATM IV, DD falls by 1.987 under Random Forest, 2.267 under GARCH, 5.850 under CRISIL and 6.087 under Altman. Reported PD does not expose these gaps because four routes are clipped to 3 bp; Altman is 4 bp higher. The IV-minus-GARCH volatility difference is $-7.775$ percentage points, so the selected short-dated ATM IV does not display a positive variance risk premium relative to the one-year GARCH measure. This comparison is qualified by horizon mismatch.

\subsection{Short-dated IV and one-year DD}
NSE stock-option maturities are below three months, while Merton DD is calculated over one year. Using short-dated ATM IV as a one-year constant volatility input assumes that current event risk and market conditions persist for the full structural horizon. If near-term volatility is temporarily low and later reverts upward, $\sigE$ and $\sigV$ are understated, so DD is biased upward and PD downward. If a near-term event inflates option volatility, the opposite bias arises. Biocon's ATM IV of 24.915\% is below the 32.690\% one-year Student-$t$ GARCH forecast, producing DD of 9.281 rather than 7.014. The 7.775 percentage-point difference may contain term-structure effects, risk-neutral pricing and liquidity effects; it should not be labelled a pure variance risk premium without maturity-matched measures. DD from IV is therefore a point-in-time scenario rather than a complete one-year forecast. (126 words)

\subsection{Did Random Forest beat GARCH, and does better volatility imply better PD?}
Random Forest did not beat GARCH on the untouched 252-day test sample. Its RMSE, MAE and QLIKE are 0.02430, 0.01811 and 0.01076, compared with 0.01990, 0.01520 and 0.00601 for GARCH; its GARCH-relative out-of-sample $R^2$ is $-0.4921$. The forest nevertheless produces a plausible main-date volatility of 31.480\%, close to GARCH's 32.690\%, and both reported PDs reach the 3 bp floor. A better volatility forecast does not automatically guarantee a better PD estimate. Structural PD also depends on market equity, the debt default point, rates, horizon and Merton's simplified capital-structure assumptions. In addition, the PD floor compresses differences in very low raw probabilities. Forecast evaluation should therefore be separated from PD interpretation: GARCH is preferred empirically, but neither route directly validates the true default probability. (124 words)

\subsection{Recommendation for a bank early-warning system}
For a bank monitoring Indian corporates, Student-$t$ GARCH-Merton DD should be the primary quantitative early-warning indicator. It updates with market prices, captures volatility clustering and heavy tails, and outperforms Random Forest on Biocon's one-year test horizon. DD should be monitored rather than only the floor-adjusted PD because DD preserves risk variation when reported PD is clipped at 3 bp. The measure should not operate alone. Agency ratings provide through-the-cycle assessment and qualitative credit analysis, while Altman EM highlights deterioration in accounting liquidity, profitability and solvency. A practical dashboard should therefore track GARCH-Merton DD frequently, use option IV as a forward-looking market check, and trigger analyst review when rating or Altman signals diverge materially. Limits are market noise, model assumptions, stale financial statements and the simplified KMV default point; judgement and exposure-specific information remain necessary. (124 words)

\subsection{Conclusion}
Biocon's market routes indicate a substantial asset cushion, conditional on observed equity value, the debt proxy and volatility. GARCH is preferred for monitoring; ATM IV, CRISIL and Altman provide complementary forward-looking, through-the-cycle and accounting checks. Random Forest adds nonlinear features but does not improve out-of-sample accuracy. The routes are complementary rather than interchangeable.

% ---------- Appendices excluded from the 10-page technical section ----------
\clearpage
\appendix
\renewcommand{\thepage}{A-\arabic{page}}
\setcounter{page}{1}

\section{Complete Model Assumption Assessment}\label{app:assumptions}
\small
\begin{longtable}{p{0.27\textwidth}p{0.27\textwidth}p{0.38\textwidth}}
\caption{BSM assumptions, Biocon assessment and consequence of failure}\\
\toprule Assumption & Biocon assessment & Consequence if violated\\\midrule\endfirsthead
\toprule Assumption & Biocon assessment & Consequence if violated\\\midrule\endhead
GBM with normal log returns and constant drift/volatility & Approximate; returns are heavy-tailed and volatility varies & BSM prices and IV are biased; smile and term structure emerge\\
Constant volatility and risk-free rate over option life & Volatility restrictive; short-rate approximation reasonable & Incorrect option discounting and unstable IV interpretation\\
Known continuous dividend yield & Discrete dividends represented through trailing yield & Forward price, bounds and IV change if dividends are mismeasured\\
Frictionless, continuous trading with unrestricted funding and short sales & Does not hold literally & Exact replication fails and prices include transaction/funding effects\\
No persistent arbitrage & Reasonable benchmark & BSM identities and bounds need not hold\\
European exercise & Appropriate for the NSE contracts used & Early exercise would require another model\\
Fair synchronous settlement and spot & Approximate; reconciled to equity close & IV can reflect stale prices or liquidity rather than expectations\\
\bottomrule
\end{longtable}

\begin{longtable}{p{0.27\textwidth}p{0.27\textwidth}p{0.38\textwidth}}
\caption{Merton assumptions, Biocon assessment and consequence of failure}\\
\toprule Assumption & Biocon assessment & Consequence if violated\\\midrule\endfirsthead
\toprule Assumption & Biocon assessment & Consequence if violated\\\midrule\endhead
Firm assets follow GBM with constant $\sigV$ & Approximate; asset value is unobservable and risk changes & Estimated $V$, $\sigV$, DD and PD may be biased\\
One zero-coupon debt claim maturing at $T$ & Biocon has multiple maturities; KMV default point is a proxy & Refinancing and maturity structure are compressed\\
Default occurs only at $T$ if $V<D$ & Restrictive; liquidity and covenants can cause earlier distress & Early default paths are ignored\\
Equity is a European call on assets & Useful structural approximation & Multiple claims and stakeholder priorities weaken the analogy\\
No bankruptcy costs, taxes, dividends or new debt before $T$ & Does not hold literally & Firm and recovery values differ from the frictionless model\\
Fixed one-year rate and sufficient equity information to infer assets & Reasonable assignment calibration but model-dependent & Results are sensitive to $E$, $D$, rate and volatility measurement\\
\bottomrule
\end{longtable}

\section{Data Sources, Timing and Cleaning Audit}\label{app:data}
\small
\begin{longtable}{p{0.23\textwidth}p{0.34\textwidth}p{0.35\textwidth}}
\caption{Principal data sources and uses}\\
\toprule Source & Data used & Treatment\\\midrule\endfirsthead
\toprule Source & Data used & Treatment\\\midrule\endhead
Biocon/NSE market workbook & Adjusted prices/returns, close, Nifty 50, India VIX, results months, shares, dividends & Adjusted series used directly; close used for equity value; exact shares applied on allotment dates\\
NSE option bhavcopy & 1,541 F\&O rows over 12 dates & Futures preserved in raw audit; options filtered sequentially\\
RBI auctions & 91-day and 364-day cut-off yields & Latest available yield by cutoff; $r=\ln(1+R/100)$\\
Biocon consolidated statements & Quarterly borrowings and FY2022--FY2026 Altman inputs & Latest public balance sheet carried forward; no growth forecast\\
CRISIL & AA+ rationale and FY2026 default study & AA category CDR, adjacent bucket and 3 bp floor\\
\bottomrule
\end{longtable}

\begin{table}[H]
\centering
\caption{Sequential option-cleaning reconciliation}
\begin{tabular}{lrr}
\toprule Stage & Removed & Remaining\\
\midrule
Raw F\&O rows & -- & 1,541\\
Remove futures & 36 & 1,505\\
Wrong expiry & 859 & 646\\
Zero volume & 199 & 447\\
Outside $0.90\leq K/S\leq1.10$ & 143 & 304\\
Lower-bound violations & 2 & 302\\
\bottomrule
\end{tabular}
\end{table}

\section{Detailed Results and Independent Validation}\label{app:results}
\small
\begin{table}[H]
\centering
\caption{Main-date cross-route structural outputs}
\begin{tabular}{lrrrr}
\toprule Route & $\sigE$ & $\sigV$ & DD & Reported PD\\
\midrule
ATM IV & 24.915\% & 21.653\% & 9.281 & 3 bp\\
Student-$t$ GARCH & 32.690\% & 28.410\% & 7.014 & 3 bp\\
Random Forest & 31.480\% & 27.359\% & 7.295 & 3 bp\\
CRISIL reverse calibration & 63.124\% & 54.862\% & 3.432 & 3 bp\\
Altman EM reverse calibration & 67.098\% & 58.319\% & 3.195 & 7 bp\\
\bottomrule
\end{tabular}
\end{table}

\begin{itemize}
\item Own IV solver: 302/302 convergence; maximum repricing error INR 0.00009187; maximum absolute difference from \texttt{py\_vollib} approximately 0.00000337.
\item Merton: every SciPy \texttt{fsolve} solution substituted back into both equations; check-only \texttt{merton} package agrees within stated tolerances.
\item GARCH: all normal and Student-$t$ rolling fits converge and satisfy stationarity; Student-$t$ BIC is lower on 12/12 dates.
\item Random Forest: ordered, purged 60/20/20 split; final test untouched during model selection; GARCH-relative $R^2$ reported honestly.
\item Rating and Altman: every reverse solve checked with an independent bounded solver.
\end{itemize}

\section{References and Reproducibility}
\begin{thebibliography}{9}
\bibitem{altman1968} Altman, E. I. (1968). Financial ratios, discriminant analysis and the prediction of corporate bankruptcy. \textit{Journal of Finance}, 23(4), 589--609.
\bibitem{merton1974} Merton, R. C. (1974). On the pricing of corporate debt: The risk structure of interest rates. \textit{Journal of Finance}, 29(2), 449--470.
\bibitem{crisil2025} CRISIL Ratings. Biocon Limited rating rationale, 16 December 2025. \url{https://www.crisil.com/mnt/winshare/Ratings/RatingList/RatingDocs/BioconLimited_December%2016_%202025_RR_385018.html}.
\bibitem{crisilstudy} CRISIL Ratings. \textit{Annual Default and Rating Transition Study FY2026}.
\bibitem{rbi} Reserve Bank of India. 91-day and 364-day Government of India Treasury-bill auction results, assignment period.
\bibitem{nse} National Stock Exchange of India. Biocon historical market data, F\&O bhavcopy and corporate-action filings.
\bibitem{biocon} Biocon Limited. Consolidated annual reports, quarterly results and share-allotment filings used in the data appendix.
\end{thebibliography}

The reproducible notebook pins \texttt{arch 8.0.0}, \texttt{statsmodels 0.15.0}, \texttt{scikit-learn 1.9.1}, \texttt{py-vollib 1.0.12}, \texttt{merton 1.0.2} and \texttt{openpyxl 3.1.5}. SciPy supplies the authoritative Merton solver; \texttt{arch} estimates GARCH; \texttt{statsmodels} supplies diagnostics; scikit-learn estimates Random Forest; \texttt{py\_vollib} and \texttt{merton} are check-only references.

\end{document}
